\documentclass[a4paper,landscape,headrule,footrule]{foils}
%%
%%% macros for Theories of Grammar
%%%
\usepackage{polyglossia}
\setdefaultlanguage{english}
%\setmainfont{TeX Gyre Pagella}


\newcommand{\logo}{~}
\newcommand{\header}[3]{%
\title{\vspace*{-2ex} \large JMORF/MORS --- Morpho-Syntax
\\[2ex] \Large  \red{#2} \\ \red{#3}}
\author{\blu{Francis Bond}   \\ 
\normalsize  \textbf{Palacký University}\\
\normalsize  \url{https://fcbond.github.io/}\\
\normalsize  \texttt{bond@ieee.org}}
\MyLogo{JMORF (2026)}
\renewcommand{\logo}{#2}
\hypersetup{
   pdfinfo={
     Author={Francis Bond},
     Title={#1: #2},
     Subject={JMORF/MORS: Morpho-Syntax},
     Keywords={Syntax, Morphology, HPSG, Unification, Constructions},
     License={CC BY 4.0}
   }
 }

\date{#1
  \\ Location: KB-1.11}
}
\usepackage{hyperref}
\hypersetup{
     colorlinks,
     linkcolor={blue!50!black},
     citecolor={red!50!black},
     urlcolor={blue!80!black}
   }



\usepackage{xcolor}
\usepackage{graphicx}
\newcommand{\blu}[1]{\textcolor{blue}{#1}}
\newcommand{\red}[1]{\textcolor{red}{#1}}
\newcommand{\grn}[1]{\textcolor{green}{#1}}
\newcommand{\hide}[1]{\textcolor{white}{#1}}
\newcommand{\emp}[1]{\textcolor{blue}{#1}}
\newcommand{\txx}[1]{\textbf{\textcolor{blue}{#1}}}
\newcommand{\lex}[1]{\textbf{\mtcitestyle{#1}}}


% \renewcommand{\labelitemi}{\textcolor{violet}{‣}}
% \renewcommand{\labelitemii}{\textcolor{purple}{‣}}

\newcommand{\subhead}[1]{\noindent\textbf{#1}\\[5mm]}

\newcommand{\Bad}{\emp{\raisebox{0.15ex}{\ensuremath{\mathbf{\otimes}}}}}
\newcommand{\bad}[1]{*\eng{#1}}

\newcommand{\com}[1]{\hfill (\emp{#1})}%

\usepackage{relsize,xspace}
\newcommand{\into}{\ensuremath{\rightarrow}\xspace}
\newcommand{\tot}{\ensuremath{\leftrightarrow}\xspace}
\usepackage{url}
\newcommand{\lurl}[1]{\MyLogo{\url{#1}}}

\usepackage{langsci-gb4e}
\let\eachwordone=\slshape
%%% langsci-gb4e defines \thexnumi as \@roman\@xsi{xnumi}, which only works
%%% inside an exe environment; some decks use xlisti on its own.  Restore
%%% the plain gb4e definition so that keeps working.
\makeatletter
\def\thexnumi{\@xsi{xnumi}}
\makeatother
\newcommand{\lx}[1]{\textbf{\mtciteform{#1}}}
\newcommand{\ix}{\ex\slshape}



\newcommand{\ent}{\ensuremath{\Rightarrow}\xspace}
\newcommand{\ngv}{\ensuremath{\not\Rightarrow}\xspace}
%\usepackage{times}
%\usepackage{nttfoilhead}
\newcommand{\myslide}[1]{\foilhead[-25mm]{\raisebox{12mm}[0mm]{\red{#1}}}\MyLogo{\logo}}
\newcommand{\myslider}[1]{\rotatefoilhead[-25mm]{\raisebox{12mm}[0mm]{\red{#1}}}}
%\newcommand{\myslider}[1]{\rotatefoilhead{\raisebox{-8mm}{\emp{#1}}}}

\newcommand{\section}[1]{\myslide{}{\begin{center}\Huge \red{#1}\end{center}}}



\usepackage[lyons,j,e,k]{mtg2e}
%\renewcommand{\mtcitestyle}[1]{\textcolor{-red!75!green!50}{\textsl{#1}}}
\renewcommand{\mtcitestyle}[1]{\textcolor{teal}{\textsl{#1}}}
\newcommand{\cs}{\mtciteform}
\newcommand{\iz}[1]{\texttt{\textup{#1}}}
\newcommand{\gm}{\textsc}
\usepackage[normalem]{ulem}
\newcommand{\ul}{\uline}
\newcommand{\ull}{\uuline}
\newcommand{\wl}{\uwave}
\newcommand{\vs}{\ensuremath{\Leftrightarrow}~}
%%%
%%% Bibliography
%%%
\usepackage{natbib}
%\usepackage{url}
\usepackage{bibentry}



\usepackage[utf8]{inputenc}
%%% trees
\usepackage{multicol}
%%% pst-node: one example in lec-12-control draws a connector between
%%% two points with \rnode/\ncbar.  This used to come in with rtrees.
\usepackage{pst-node}
\usepackage{forest}
\useforestlibrary{linguistics}
\forestapplylibrarydefaults{linguistics}

%%% \makebox's optional arguments are read as AVM delimiters by
%%% langsci-avm, so hide them behind a macro with no visible brackets.
\newcommand{\lbox}[1]{\makebox[0mm][l]{#1}}

%%% Tree styles.
%%%   words  -- trees whose leaves are actual words (italicised, aligned)
%%%   cat    -- trees whose leaves are categories or feature structures
\forestset{
  words/.style={for tree={s sep=7mm, l sep=5mm,
                          if n children=0{font=\itshape, tier=terminal}{}}},
  cat/.style={for tree={s sep=7mm, l sep=5mm}},
}
\newcommand{\wtree}[1]{\begin{forest} words #1 \end{forest}}
\newcommand{\ctree}[1]{\begin{forest} cat #1 \end{forest}}
%%% the rule being applied at this step of a derivation
\newcommand{\step}[1]{\begin{center}\emp{#1}\end{center}}

%%% AVMs (Language Science Press)
\usepackage{langsci-avm}

%%% Macros the decks use, carried over from headx.tex
\newcommand{\blank}{\rule{3em}{1pt}\xspace}
\newcommand{\ft}[1]{\textsc{#1}}
\newcommand{\idn}{\mtciteform}    %%% Indonesian
\newcommand{\prd}[1]{\textbf{#1}}
\newcommand{\typ}[1]{\textit{#1}}
%%% \val: provide before renewing, in case a package defines it.
\providecommand{\val}[1]{}
\renewcommand{\val}[1]{\textit{#1}}
%%% Manning's \avmbox, still used in tree labels and in maths where a
%%% bare langsci \tag would not work.  \avm{\tag{...}} gives the same
%%% boxed tag as the ones inside feature structures.
\newcommand{\avmbox}[1]{\mbox{\avm{\tag{#1}}}}

%%% Ring a value in a dotted oval.  These used pstricks \ovalnode,
%%% which came in with rtrees; tikz (via forest) does the same job.
\newcommand{\OV}[1]{\tikz[baseline]{\node[draw,dotted,red,ellipse,
          inner sep=1pt,outer sep=0pt]{#1};}}
\newcommand{\OVB}[1]{\tikz[baseline]{\node[draw,dotted,blue,ellipse,
          inner sep=1pt,outer sep=0pt]{#1};}}

%%% From CSLI book
\newcommand{\mc}{\multicolumn}
\newcommand{\HD}{\textbf{H}\xspace}
\newcommand{\el}{< >}       %%% empty list; langsci-avm needs the space

\newcommand{\A}{\noindent\textbf{A}: }
\newcommand{\Q}{\noindent\textbf{Q}: }

%%% Question and answer on two pages, written once:
%%%   \qaslide{Title}{... question ... \answer{... answer ...}}
%%% prints the slide twice, first without the \answer parts and then
%%% with them, so the answer appears "on click".  The example counter
%%% is reset so both pages show the same example numbers: gb4e keeps the
%%% running number in the equation counter between exe environments, and
%%% in xnumi inside one.
\newif\ifshowanswer
\newcommand{\answer}[1]{\ifshowanswer#1\fi}
\newcounter{qaxnumi}
\newcounter{qaeqn}
\newcommand{\qaslide}[2]{%
  \setcounter{qaxnumi}{\value{xnumi}}%
  \setcounter{qaeqn}{\value{equation}}%
  \showanswerfalse\myslide{#1}#2%
  \setcounter{xnumi}{\value{qaxnumi}}%
  \setcounter{equation}{\value{qaeqn}}%
  \showanswertrue\myslide{#1}#2%
  \showanswerfalse}
%\newcommand{\C}{\noindent\textbf{C}: }



%\usepackage{times}
\usepackage{bchart}
%\usepackage{pgfplots}
\begin{document}
\header{Lecture 9}{Non-referential NPs, Expletives, and Extraposition}{}
\maketitle




\myslide{Overview}
\MyLogo{\citet{sag:wasow:bender:2003} --- Chapter 11}

\begin{itemize}
\item Review: the copula \lex{be}
\item Existentials (\eng{There is an X}, \ldots)
\item Dummy \lex{it}
\item Complementizers and CP subjects
\item Extraposition (\eng{It worries me that X}, \ldots)
\item Idioms (\eng{X takes advantage of Y}, \ldots)
\item Problem Set 9: SWB 11.1, 11.3, 11.4
\end{itemize}

\section{Review}

\myslide{Introduction}
\MyLogo{\citet{sag:wasow:bender:2003} --- Chapter 11}

\begin{itemize}
\item In Chapter~10, we met the passive \lex{be}.
\item Passive \lex{be} is just a special case --- that \lex{be}
generally introduces [\ft{pred} +] constituents.
%(next slide).
\item Today, we’ll introduce another \lex{be}, which 
occurs in existential sentences starting with 
\lex{there}, e.g.\ \eng{There is a monster in Loch Ness.}
\item Then we’ll look at this use of \lex{there}.
\item Which will lead us to a more general 
examination of NPs that don’t refer, including 
some uses of \lex{it} and certain idiomatic uses of 
NPs.
\end{itemize}



\myslide{Chapter 10 Entry for \lex{be}}

\avm{< \textnormal{be},\ [\type*{be-lxm}\\
  arg-st & < {\1} , [syn & [head & [\type*{verb}\\
  form &  pass ]\\
  val & [spr & < {\1} >\\
  comps & < > ]]\\
  sem & [index & $s$ \ ]] > \\ %{\3}
  sem & [index & \ $s$\\ %{\3}\\
  restr & {< >} ] ] >
}

This takes only passive lexemes: \avm{[ form &  pass ]
}.

\myslide{Copula (Generalized)}
\MyLogo{\citet[p.~334]{sag:wasow:bender:2003}}
\avm{< \textnormal{be},\ [\type*{be-lxm}\\
  arg-st & < {\1} , [syn & [head & [ pred & $+$ ]\\
  val & [spr & < {\1} > \\ 
  comps & < > ]]\\
  sem & [index & $s$ \ ]] > \\ %{\3}
  sem & [index & \ $s$\\ %{\3}\\
  restr & {< >} ] ] >
}
\begin{center}
  \eng{Kim is happy.}\qquad \eng{Kim is sleeping.}
\end{center}
\vspace{-2ex}
\begin{itemize}\addtolength{\itemsep}{-2ex}
\item Allow also (\txx{pred}icative) adjectives \lex{fond}, \lex{sleeping}
  are [\ft{pred} +]
\item \lex{mere} is [\ft{pred} $-$] %\begin{avm}\[ pred &  - \]\end{avm}.
\item Most adjectives can be either  
\end{itemize}


\section{Existentials}

\qaslide{Existentials}{

\begin{itemize}\addtolength{\itemsep}{-1ex}
\item The \lex{be} in \eng{There is a page missing} cannot be the 
same \lex{be} that occurs in sentences like \eng{Pat is tall} or 
\eng{A cat was chased by a dog}.  Why not?
  \answer{\\ {\small\emp{Here \lex{be} has two complements, and \lex{there} is not the subject of \eng{missing}.  Nor is \eng{a page missing} one [\ft{pred} +] phrase: it is not even a constituent (\bad{A page missing caused problems}).}}}
\item So we need a separate lexical entry for this \lex{be}, 
  stipulating:
  \begin{itemize}
  \item Its \ft{spr} must be \lex{there}.
  \item It takes two complements, the first an NP and the 
    second an AP, PP, or (certain kind of) VP.
  \item The semantics should capture the relation between, e.g. 
    \eng{There is a page missing} and \eng{A page is missing}.  
  \end{itemize}
\item More examples: \eng{There was a person at the door};
  \eng{There are many people fond of linguistics};
  \eng{There are people looking at us}.
\end{itemize}
}

\myslide{Lexical Entry for the Existential \lex{be}}

\scalebox{0.95}{\avm{< \textnormal{be},\ [\lbox{\textit{exist-be-lxm}}\\
  arg-st & < !\shortstack{NP\\\avm{[ form
    & there ]}}!, {\1},
  [syn & [head & [ pred & $+$ ]\\
  val & [spr & < {\1} > \\ comps & < > ]]\\
  sem & [index & $s$ \ ]] > \\ %{\3}
  sem & [index & \ $s$\\ %{\3}\\
  restr & {< >} ] ] >
}}


\qaslide{Questions About the Existential \lex{be}}{

\begin{itemize}
\item What type of constituent is the third argument?
  \answer{\\ {\small\emp{A predicative phrase (AP, PP, or VP) still missing its subject: \ft{spr} \avm{< \1 >}, \ft{comps} \avm{< >}.}}}
\item Why is the third argument [\ft{pred} +]?
  \answer{\\ {\small\emp{Only predicative phrases go there: \eng{There is a page missing}, not \bad{There is a page lose}.}}}
\item Why is the second argument tagged as identical to the \ft{spr}
  of the third argument?
  \answer{\\ {\small\emp{The NP is understood as the subject of the predicate: \eng{a page} is what is \eng{missing}.}}}
\end{itemize}

\begin{center}
  \begin{tiny}
    \avm{< \textnormal{be},\ [\lbox{\textit{exist-be-lxm}}\\
      arg-st & < !\shortstack{NP\\\avm{[ form
          & there]}}!, {\1},
      [syn & [head & [ pred & $+$ ]\\
      val & [spr & < {\1} > \\ comps & < > ]]\\
      sem & [index & $s$ \ ]] > \\ %{\3}
      sem & [index & \ $s$\\ %{\3}\\
      restr & {< >} ] ] >
}
  \end{tiny}
\end{center}
}

\qaslide{There Are Questions Left}{

\begin{itemize}\addtolength{\itemsep}{-1.5ex}
\item What is the contribution of this \lex{be} to the semantics of the sentences 
it occurs in?
  \answer{\\ {\small\emp{None: its \ft{restr} is empty and its \ft{index} is the third argument's.  The meaning comes from the predicate.}}}
\item Can all [\ft{pred} +] predicates appear as the third argument in 
existentials?
  \answer{\\ {\small\emp{No: predicative NPs are [\ft{pred} +] too, and so are adjectives like \lex{tall} (\bad{There is a man tall}).}}}
\item How do we rule out *\lex{There was a greyhound a good runner}?
  \answer{\\ {\small\emp{The entry as it stands doesn't: we would need to add a constraint, e.g.\ that the third argument's \ft{head} is not \val{noun}.}}}  
\end{itemize}

\begin{center}
  \begin{tiny}
    \avm{< \textnormal{be},\ [\lbox{\textit{exist-be-lxm}}\\
      arg-st & < !\shortstack{NP\\\avm{[ form
          & there ]}}!, {\1},
      [syn & [head & [ pred & $+$ ]\\
      val & [spr & < {\1} > \\ comps & < > ]]\\
      sem & [index & $s$ \ ]] > \\ %{\3}
      sem & [index & \ $s$\\ %{\3}\\
      restr & {< >} ] ] >
}
  \end{tiny}
\end{center}
}

\myslide{The Entry for Existential \lex{there}}

\begin{center}
  \avm{< \textnormal{there},\ [\lbox{\textit{pron-lxm}}\\
    syn & [head & [ form & there \\
    agr  & [ per & 3rd ]   ] ]\\
    sem & [mode & none \\
    index & none \\
    restr & < > ]
    ] >
}
\end{center}

\qaslide{Questions About Existential \lex{there}}{

\begin{itemize}\addtolength{\itemsep}{-1.5ex}
\item Why do we call it a pronoun?
  \answer{\\ {\small\emp{It is a whole NP by itself, and can be a subject (\eng{Is there a page missing?}).}}}
\item Why don't we give it a value for \ft{num}?
  \answer{\\ {\small\emp{The verb agrees with the NP after \lex{be}: \eng{There is a cat}, \eng{There are cats} (P1 of the problem set).}}}
\item What does this entry claim is \lex{there}’s contribution to the
  semantics of the sentences it appears in?
  \answer{\\ {\small\emp{Nothing: \ft{mode} and \ft{index} \val{none}, \ft{restr} empty.}}}
\item Is this a correct claim?
  \answer{\\ {\small\emp{Mostly: it doesn't refer, and can't be questioned (\bad{What is a page missing?}).}}}
\end{itemize}

\begin{center} \tiny
  \avm{< \textnormal{there},\ [\lbox{\textit{pron-lxm}}\\
    syn & [head & [ form & there \\
    agr  & [ per & 3rd ]   ] ]\\
    sem & [mode & none \\
    index & none \\
    restr & < > ]
    ] >
}
\end{center}
}

\myslide{Sample Tree for Existential \lex{be}}

\begin{exe}
  \ex \eng{There are cats in boxes.}
\end{exe}

\begin{minipage}[t]{0.45\textwidth}
\begin{forest}
cat
[{S} [{NP} [{There}]]
[{VP} [{V} [{are}]]
[{NP} [{N} [{cats}]]]
[{PP} [{P} [{in}]]
[{NP} [{N} [{boxes}]]]]]]
\end{forest}
\end{minipage}\hfill
\begin{minipage}[t]{0.5\textwidth}
\begin{itemize}
\item \eng{cats} and \eng{in boxes} are sisters: two complements of
  \lex{are}.
\item \eng{cats in boxes} is not a constituent here (compare
  \eng{Cats in boxes are cute}, where it is).
\end{itemize}
\end{minipage}

\myslide{Step by Step: What \lex{are} Asks For}

\begin{center}
\scalebox{0.85}{\begin{forest}
cat, for tree={align=center}
[{S}
[{\avmbox{1}NP\\ {\small$[$ \ft{form} there $]$}} [{There}]]
[{VP\\ {\small$[$ \ft{spr} $\langle \avmbox{1} \rangle$, \ft{comps} $\langle\ \rangle$ $]$}}
[{V\\ {\small$[$ \ft{spr} $\langle \avmbox{1} \rangle$,}\\ {\small\ft{comps} $\langle \avmbox{2}, \avmbox{3} \rangle$ $]$}} [{are}]]
[{\avmbox{2}NP$_i$} [{cats}, roof]]
[{\avmbox{3}PP\\ {\small$[$ \ft{pred} +, \ft{spr} $\langle \avmbox{2} \rangle$ $]$}} [{in boxes}, roof]]]]
\end{forest}}
\end{center}
\vspace{-2ex}
\begin{itemize}\addtolength{\itemsep}{-1.5ex}
\item The \ft{arg-st} of existential \lex{be} is $\langle$ NP[\ft{form}
  there], \avmbox{2}, \avmbox{3} $\rangle$; the ARP puts the first on
  \ft{spr} and the rest on \ft{comps}.
\item The PP \eng{in boxes} is [\ft{pred} +] and still wants a subject
  (\ft{spr} $\langle \avmbox{2} \rangle$): that subject is \eng{cats},
  so the cats are in the boxes.
\end{itemize}

\qaslide{Which Rule Licenses Each Node?}{
\begin{minipage}[t]{0.4\textwidth}
\begin{forest}
cat
[{S} [{NP} [{There}]]
[{VP} [{V} [{are}]]
[{NP} [{N} [{cats}]]]
[{PP} [{P} [{in}]]
[{NP} [{N} [{boxes}]]]]]]
\end{forest}
\end{minipage}\hfill
\begin{minipage}[t]{0.55\textwidth}
\begin{itemize}\addtolength{\itemsep}{-1ex}
\item S?
  \answer{\\ {\small\emp{Head-Specifier Rule: the VP and its specifier \eng{There}.}}}
\item VP?
  \answer{\\ {\small\emp{Head-Complement Rule: \eng{are} and both its complements, NP and PP.}}}
\item PP?
  \answer{\\ {\small\emp{Head-Complement Rule: \eng{in} and its NP.}}}
\item NP over \eng{cats}, NP over \eng{boxes}?
  \answer{\\ {\small\emp{Head-Complement Rule with no complements (a plural noun's specifier is optional).}}}
\item NP over \eng{There}?
  \answer{\\ {\small\emp{None needed: a pronoun is already a saturated NP.}}}
\end{itemize}
\end{minipage}
}

\qaslide{Which Lexical Rule Licenses Each Word?}{
\begin{itemize}\addtolength{\itemsep}{-1ex}
\item \eng{There}?
  \answer{\\ {\small\emp{Constant Lexeme LR: \val{pron-lxm} is a \val{const-lxm}; the word is the lexeme unchanged.}}}
\item \eng{are}?
  \answer{\\ {\small\emp{Non-3rd-Singular Verb LR applied to existential \lex{be} (\val{exist-be-lxm}); the form is irregular, \eng{are}.  Agreement with \eng{cats} is Problem Set 9, P1.}}}
\item \eng{cats}, \eng{boxes}?
  \answer{\\ {\small\emp{Plural Noun LR (\val{cntn-lxm} \into\ word, add \textit{-s}).}}}
\item \eng{in}?
  \answer{\\ {\small\emp{Constant Lexeme LR, from the predicative preposition lexeme (\val{predp-lxm}, [\ft{pred} +]).}}}
\end{itemize}
}

\section{Dummy \lex{it}}

\myslide{Other NPs That Don’t Seem to Refer}

{\large
\begin{exe}
\ex \eng{\ul{It} sucks that the Rockies lost the series.}
\ex \eng{\ul{It} is raining.}
\ex \eng{Andy took \ul{advantage} of the opportunity.}
\ex \eng{Lou kicked \ul{the bucket}.}
\end{exe}
}

\myslide{What About \eng{It follows that you are wrong}?}
\MyLogo{$^*$We need these anyway (independently motivated)}
\begin{itemize}
\item This is an example of \txx{extraposition}.
\item To analyze it, we need:
  \begin{itemize}
  \item An analysis of this use of \lex{that}$^*$
  \item Entries for verbs that take clausal subjects$^*$
    \begin{exe}
      \ex \label{s:that:follows} \eng{\ul{That you are wrong} follows.}
    \end{exe}
  \item A lexical entry for dummy \lex{it}
  \item A rule to account for the relationship 
    between pairs like  (\ref{s:that:follows}) and (\ref{s:follows:that})
\begin{exe}
  \ex \label{s:follows:that}
  \eng{\ul{It} follows \ul{that you are wrong}.}
\end{exe}
\end{itemize}
\end{itemize}

\myslide{The Entry for Dummy \lex{it}}

\begin{center}
  \avm{< \textnormal{it},\ [\lbox{\textit{pron-lxm}}\\
    syn & [head & [ form & it \\
    agr  & 3sing   ] ]\\
    sem & [mode & none \\
    index & none \\
    restr & < > ]
    ] >
}
\end{center}

\qaslide{Questions About Dummy \lex{it}}{


\begin{itemize}
\item How does it differ from the entry for dummy \lex{there}?
Why do they differ in this way?
  \answer{\\ {\small\emp{\ft{form} \val{it}, and \ft{agr} \val{3sing}, not just 3rd person: dummy \lex{it} always takes a singular verb (\eng{It is raining}), while \lex{there} agrees with the NP after it.}}}
\item Is this the only entry for \lex{it}?
  \answer{\\ {\small\emp{No: referential \lex{it} (\eng{I saw it}) is a separate entry, with \ft{mode} \val{ref} and an index.}}}
\end{itemize}
\begin{center}\small
  \avm{< \textnormal{it},\ [\lbox{\textit{pron-lxm}}\\
    syn & [head & [ form & it \\
    agr  & 3sing   ] ]\\
    sem & [mode & none \\
    index & none \\
    restr & < > ]
    ] >
}
\end{center}
}

\section{Complementizers and CP Subjects}

\myslide{A New Type of Lexeme: Complementizers}

\begin{center}\large
  \avm{\textit{comp-lxm}:  [ \\
    syn & [head & [\lbox{\textit{comp}} \\
    agr  & 3sing   ] \\
    val & [ spr & < > ] ]\\
    arg-st & < !\shortstack{S \\ \avm{[ index & s ]}}!>\\
    sem & [ index & s \\
    restr & < > ]
    ]
}
\end{center}


\qaslide{Questions About the Type \val{comp-lxm}}{

\begin{itemize}\addtolength{\itemsep}{-1.5ex}
\item Why does it stipulate values for both \ft{spr} and \ft{arg-st}?
  \answer{\\ {\small\emp{With \ft{spr} \avm{< >}, the ARP must put the S on \ft{comps}: \lex{that} takes it as a complement, not a subject.}}}
\item Why is its \ft{index} value the same as its argument’s?
  \answer{\\ {\small\emp{A CP is about the same situation as its S, so a verb that takes a CP gets the S's situation.}}}
\item What is its semantic contribution?
  \answer{\\ {\small\emp{None: its \ft{restr} is empty.  (Each complementizer adds a \ft{mode}, e.g.\ \val{prop} for \lex{that}.)}}}
\end{itemize}
\begin{center}\small
  \avm{\textit{comp-lxm}:  [ \\
    syn & [head & [\lbox{\textit{comp}} \\
    agr  & 3sing   ] \\
    val & [ spr & < > ] ]\\
    arg-st & < !\shortstack{S \\ \avm{[ index & s ]}}!>\\
    sem & [ index & s \\
    restr & < > ]
    ]
}
\end{center}
}

\myslide{The Type \val{comp}}


\begin{center}
  \begin{forest}
cat, for tree={font=\itshape}
[{!\shortstack{pos\\ \avm{[ \textsc{form, pred} ]}}!} [{agr-pos \avm{[ \textsc{agr} ]}} [{verb \avm{[ \textsc{aux} ]}}]
[{nominal \avm{[ \textsc{case} ]}} [{\shortstack{noun \\ \mbox{} }}]
[{!\shortstack{comp \\\avm{[ \textsc{form} & cform ]}}!}]]
[{det \avm{[ \textsc{count} ]}}]]
[{prep}]
[{adj}]
[{conj}]]
\end{forest}
\end{center}

\myslide{The Lexical Entry for Complementizer \lex{that}}
\bigskip
\begin{center}\large
  \avm{< \textnormal{that},\ [\lbox{\textit{comp-lxm}}\\
    arg-st  & < [ form & fin ]>\\
    sem & [mode & prop ]
    ] >
}
\end{center}
\qaslide{\ldots with Inherited Information Filled In}{

\begin{center}\small
  \avm{< \textnormal{that},\ [\lbox{\textit{comp-lxm}}\\
   syn & [head & [\lbox{\textit{comp}} \\
    form & \OV{cform} \\
    agr  & 3sing   ] \\
    val & [ spr & < > ] ]\\
    arg-st  & < !\shortstack{S \\ \avm{[ index & s \\ form & fin]} }!>\\
    sem & [mode & prop \\
    index & s \\
    restr & < > ]
    ] >
}
\end{center}

\begin{itemize}
\item Question: Where did [\ft{form} \val{cform}] come from?
  \answer{\\ {\small\emp{From the head type \val{comp}, which says [\ft{form} \val{cform}] (see the \val{pos} hierarchy).}}}
\end{itemize}
}

\myslide{Structure of a Complementizer Phrase}

\begin{exe}
  \ex \eng{that the Giants lost}
\end{exe}

% \begin{center}\small
%   \begin{tree}%\renewcommand{\lf}[1]{\br{#1}{}}\avmfont{\sc}
% \br{\textnormal{CP} \begin{avm}\[ \ft{head} & \@1 \\
%           \ft{val} & \< \ft{spr} & < > \\ \ft{comps} & < > \> \]\end{avm}}%
% {\br{\textnormal{C} \begin{avm}\[ \ft{head} & \@1 \[\textit{comp} 
% \\  \ft{form} & cform \] \\
%           \ft{val} & \< \ft{spr} & < > \\ \ft{comps} & \< \@2 \> \> \]
% \end{avm}}{\lf{that}}
% \br{\@2 \textnormal{S}}{\tlf{the Giants lost}}}
% \end{tree}
% \end{center}
% \newpage
\begin{center}
  \begin{forest}
cat, for tree={font=\itshape}
[{\textnormal{CP} \avm{[ \ft{head} & \1 \\
          \ft{val} & [ \ft{spr} & < > \\ \ft{comps} & < > ] ]}} [{\textnormal{C} \avm{[ \ft{head} & \1 [\textit{comp} \\  \ft{form} & cform ] \\
          \ft{val} & [ \ft{spr} & < \mbox{} > \\ \ft{comps} & < \2 > ] ]}} [{that}]]
[{\avm{\2} \textnormal{S}} [{the Giants lost}, roof]]]
\end{forest}
\end{center}



\qaslide{Is \lex{that} Really the Head?}{
\begin{itemize}\addtolength{\itemsep}{-1ex}
\item We treat the complementizer as the head of \eng{that the Giants
    lost}, although the S carries all the meaning.  Does that seem
  strange?  What is the evidence for it?
  \answer{
  \begin{itemize}\addtolength{\itemsep}{-0.5ex}
  \item \emp{The little word decides where the phrase can go: verbs select
    a CP, and the S on its own will not do (\eng{That they won matters}
    vs.\ \bad{They won matters}).}
  \item \emp{The complementizer decides what kind of S it takes: \lex{that}
    takes a finite S, \lex{for} takes an infinitive (\eng{For Kim to
      leave would be sad}).  Selection is what heads do.}
  \item \emp{It is the same pattern as the preposition in \eng{of the
      election}: a small head, a big complement, no meaning of its own,
    the \ft{index} passed up.}
  \item \emp{The phrase is a CP, not an S, and [\ft{form} \val{cform}] is
    what \lex{know}, \lex{matter}, \ldots\ look for.}
  \end{itemize}}
\end{itemize}
}

\myslide{Sample Verb with a CP Subject}


\bigskip
\begin{center}\large
  \avm{< \textnormal{matter},\ [\lbox{\textit{siv-lxm}}\\
    arg-st  & < [ sem & [ index & \1 ] ]>\\
    sem & [index & s \\
    restr & < [ reln & \prd{matter} \\
               sit & s \\
               mattering & \1 ] > ]
    ] >
}
\end{center}

Note: the only constraint on the first argument is semantic.

\myslide{A Problem}

\begin{itemize}
\item We constrained the subject of \lex{matter} only semantically.  However\ldots{}
\item CP and S are semantically identical, but we get:
  \begin{exe} 
    \ex \eng{That they won matters} vs.\ \bad{They won matters}
  \end{exe}
\item Argument-marking PPs are semantically identical to their object 
  NPs, but we get:
  \begin{exe} 
    \ex \eng{The election mattered} vs.\ \bad{Of the election mattered}
  \end{exe}
\item So we need to add a syntactic constraint.
\end{itemize}

\newpage

\begin{center}\small
  \avm{< \textnormal{matter},\ [\lbox{\textit{siv-lxm}}\\
    arg-st  & < [ syn & [ head & \OV{nominal} ] \\
                    sem & [ index & \1 ] ]>\\
    sem & [index & s \\
    restr & <  [ reln & \prd{matter} \\
                  sit & s \\
                  mattering & \1 ] > ]
    ] >
}
\end{center}


\begin{itemize}
\item  S and PP subjects are generally impossible, so this constraint should
   probably be on \val{verb-lxm}.
 \end{itemize}
 

\section{Extraposition}

\myslide{Extraposition (at Last)}
%\MyLogo{Examples adapted from wikipedia}
\begin{itemize}
\item Extraposition alters word order so that a relatively ``heavy''
  constituent appears to the right of its canonical position.
\end{itemize}

\begin{exe}\large
  \ex 
  \begin{xlist}
    \ex  \eng{\ul{That you were wrong} follows.}
    \ex \eng{\ul{It} follows \ul{that you were wrong}.}
  \end{xlist}
  \ex 
  \begin{xlist}
    \ex  \eng{\ul{That I mistyped it} was frustrating.}
    \ex \eng{\ul{It} was frustrating \ul{that I mistyped it}.}
  \end{xlist}
  \ex 
  \begin{xlist}
   \ex \eng{Did \ul{that this happened} surprise you?}
   \ex  \eng{Did it surprise you \ul{that this happened}?}
  \end{xlist}
\end{exe}



% [1]
%   Extraposing a constituent results in a discontinuity and in this
%   regard, it is unlike shifting, which does not generate a
%   discontinuity. The extraposed constituent is separated from its
%   governor by one or more words that dominate its governor. Two types
%   of extraposition are acknowledged in theoretical syntax: standard
%   cases where extraposition is optional and it-extraposition where
%   extraposition is obligatory. Extraposition is motivated in part by a
%   desire to reduce center embedding by increasing right-branching and
%   thus easing processing, center-embedded structures being more
%   difficult to process. Extraposition occurs frequently in English and
%   related languages.


\myslide{Post-Inflectional Lexical Rule: \val{pi-rule}}
\begin{center}%\small
  \avm{\val{pi-rule}: [\type*{l-rule}
        input & < / \0, [\type*{word}
            syn & [ head & / \1 \\
            val & [ mod &  \tag{A}  ] ] ] >\\
        output & < / \0,  [\type*{word}
            syn & [ head & / \1 \\
            val & [ mod &  \tag{A}  ] ] ]  >  ]
}
\end{center}

We use this when we expect the structure to be largely formed, and we
want to change the order of things, \ldots


\qaslide{The Extraposition Lexical Rule}{

\begin{center}\small
  \avm{[\type*{pi-rule}
      input & < X, [ syn & [ val & [ spr & < \1 CP > \\ comps & \tag{A} ] ]  ]  >\\
      output & <Y,  [ syn & [val & [ spr & < NP [ form & it ] > \\
      comps & \tag{A} $\oplus$ < \1 > ] ] ] > ]
}
\end{center}

\begin{itemize}
\item Why is the type \val{pi-rule}?
  \answer{\\ {\small\emp{It takes a word to a word, after inflection: the verb has already agreed with its CP subject, and \lex{it} is 3sing too.}}}
\item Why doesn’t it say anything about the semantics?
  \answer{\\ {\small\emp{The meaning doesn't change, and \val{l-rule}s copy \ft{sem} from input to output.}}}
\item Why is the \ft{comps} value \avmbox{A} not $\langle \rangle$?
  \answer{\\ {\small\emp{Some verbs with CP subjects also have complements (next slide).  Whatever was on \ft{comps} stays, and the CP goes on the end, after them.}}}
\end{itemize}
}

\myslide{Extraposition with Verbs Whose \ft{comps} Lists Are Nonempty}


\begin{exe}\large
\ex \eng{It worries me that war is imminent.}
\ex \eng{It occurred to Pat that Chris knew the answer.}
\ex \eng{It endeared you to Andy that you wore a funny hat.}
\end{exe}

\myslide{Sample Tree with Extraposition}

\begin{exe}
  \ex \eng{It matters that Kim arrived.}
\end{exe}

\begin{center}
\scalebox{0.8}{\begin{forest}
cat
[{S} [{NP} [{It}]]
[{VP} [{V} [{matters}]]
[{CP} [{C} [{that}]]
[{S} [{NP} [{Kim}]]
[{VP} [{V} [{arrived}]]]]]]]
\end{forest}}
\end{center}

\myslide{Step by Step: What \lex{matters} Asks For}

\begin{center}
\scalebox{0.85}{\begin{forest}
cat, for tree={align=center}
[{S}
[{\avmbox{1}NP\\ {\small$[$ \ft{form} it $]$}} [{It}]]
[{VP\\ {\small$[$ \ft{spr} $\langle \avmbox{1} \rangle$, \ft{comps} $\langle\ \rangle$ $]$}}
[{V\\ {\small$[$ \ft{spr} $\langle \avmbox{1} \rangle$, \ft{comps} $\langle \avmbox{2} \rangle$ $]$}} [{matters}]]
[{\avmbox{2}CP\\ {\small$[$ \ft{index} $s$ $]$}} [{that Kim arrived}, roof]]]]
\end{forest}}
\end{center}
\vspace{-2ex}
\begin{itemize}\addtolength{\itemsep}{-1.5ex}
\item Before extraposition \lex{matters} is [\ft{spr} $\langle$ \avmbox{2}CP
  $\rangle$, \ft{comps} \el]: \eng{That Kim arrived matters}.
\item The Extraposition LR swaps in a dummy subject: [\ft{spr} $\langle$
  NP[\ft{form} it] $\rangle$, \ft{comps} $\langle$ \avmbox{2} $\rangle$].
\item The semantics is untouched: the \ft{mattering} is still the
  CP's \ft{index} $s$, the situation of Kim arriving.
\end{itemize}

\qaslide{Which Rule Licenses Each Node?}{
\begin{minipage}[t]{0.4\textwidth}
\begin{forest}
cat
[{S} [{NP} [{It}]]
[{VP} [{V} [{matters}]]
[{CP} [{C} [{that}]]
[{S} [{NP} [{Kim}]]
[{VP} [{V} [{arrived}]]]]]]]
\end{forest}
\end{minipage}\hfill
\begin{minipage}[t]{0.55\textwidth}
\begin{itemize}\addtolength{\itemsep}{-1ex}
\item Top S?
  \answer{\\ {\small\emp{Head-Specifier Rule: the VP and its specifier \eng{It}.}}}
\item VP \eng{matters that Kim arrived}?
  \answer{\\ {\small\emp{Head-Complement Rule: after extraposition the CP is a complement of \eng{matters}.}}}
\item CP?
  \answer{\\ {\small\emp{Head-Complement Rule: \eng{that} and its S complement.}}}
\item Lower S, VP \eng{arrived}?
  \answer{\\ {\small\emp{Head-Specifier Rule; Head-Complement Rule with no complements.}}}
\item NP over \eng{It}, NP over \eng{Kim}?
  \answer{\\ {\small\emp{None needed: pronouns and proper nouns are already saturated, so the word is an NP.}}}
\end{itemize}
\end{minipage}
}

\qaslide{Which Lexical Rule Licenses Each Word?}{
\begin{itemize}\addtolength{\itemsep}{-1ex}
\item \eng{It}?
  \answer{\\ {\small\emp{Constant Lexeme LR, from the dummy \lex{it} entry ([\ft{form} it], \ft{index} \val{none}).}}}
\item \eng{matters}?
  \answer{\\ {\small\emp{Two rules: the 3rd-Singular Verb LR (lexeme \into\ word, add \textit{-s}), then the Extraposition LR (word \into\ word, a \val{pi-rule}).}}}
\item \eng{that}?
  \answer{\\ {\small\emp{Constant Lexeme LR: \val{comp-lxm} is a \val{const-lxm}.}}}
\item \eng{Kim}?
  \answer{\\ {\small\emp{Constant Lexeme LR (\val{pn-lxm}).}}}
\item \eng{arrived}?
  \answer{\\ {\small\emp{Past-Tense Verb LR.}}}
\end{itemize}
}

\section{Idioms}

\myslide{Another Nonreferential Noun: \lex{advantage}}
\bigskip
\begin{center}
  \avm{< \textnormal{advantage},\ [\lbox{\textit{massn-lxm}}\\
    syn & [head & [ form & advantage \\
    agr  & 3sing   ] ]\\
    sem & [mode & none \\
    index & none \\
    restr & < > ]
    ] >
}
\end{center}


\myslide{The Verb That Selects \lex{advantage}}
\MyLogo{\textbf{take\_advantage} $\approx$ \textbf{exploit}}
\begin{center}
  \avm{< \textnormal{take},\ [\lbox{\textit{ptv-lxm}}\\
    arg-st  & < NP$_i$, [ form & advantage ], 
    [ form & of  \\
       index & j ]>\\
    sem & [index & s \\
    restr & <  [ reln & \prd{take\_advantage} \\
               sit & s \\
               exploiter & i \\
               exploited & j ] > ]
    ] >
}
\end{center}


\qaslide{Our Analyses of Idioms and Passives Interact\ldots{}}{


\begin{itemize}
\item We generate
\begin{exe} 
  \ex \eng{Advantage was taken of the situation by many people.}
  \ex \eng{Tabs are kept on foreign students.}
\end{exe}
with roughly the same meanings as
\begin{exe} 
  \ex \eng{Many people took advantage of the situation.}
  \ex \eng{Someone keeps tabs on foreign students.}
\end{exe}


\item But not:
  \begin{exe} 
    \ex \eng{Many people were taken advantage of.}
  \end{exe}
\item Why not?
  \answer{ {\small\emp{\eng{taken} is the output of the Passive LR, which makes the second \ft{arg-st} element (\eng{advantage}) the subject; \eng{many people} is inside the PP, third on the list.  This ``pseudo-passive'' would need a different rule, one that makes the PP's object the subject and leaves everything else in place.}}}
\end{itemize}
}

\myslide{Overview}
\MyLogo{\citet{sag:wasow:bender:2003} --- Chapter 11}

\begin{itemize}
\item Review: the copula \lex{be}
\item Existentials (\lex{there}, \lex{be})
\item Dummy \lex{it}
\item Complementizers and CP subjects (\lex{that})
\item Extraposition (\lex{it}, LR)
\item Idioms (\lex{take\_advantage}, \ldots)
\item Problem Set 9: SWB 11.1, 11.3, 11.4
\end{itemize}

\section{Problem Set 9}

\myslide{P1: \lex{there} and Agreement}
\MyLogo{Based on \citet[Chapter 11, Problem 1]{sag:wasow:bender:2003}}

The analysis of existential \lex{there} sentences presented so far
says nothing about verb agreement.

\begin{itemize}
\item[A.] Consult your intuitions (and/or those of your friends, if
  you wish) to determine what the facts are regarding number
 agreement of the verb in \lex{there} sentences.

 Consider these sentences: which ones do you think are ok?

  \begin{exe}
    \ex\label{s:is-sg} There is a cat on the mat.
    \ex\label{s:are-sg} There are a cat on the mat.
    \ex\label{s:is-pl} There is cats on the mat.
    \ex\label{s:are-pl} There are cats on the mat.
  \end{exe}

\newpage
  Give an informal
  statement of a generalization covering these facts, and illustrate
  it with both grammatical and ungrammatical examples.  [{\sl Note:
    Intuitions vary regarding this question, across both individuals
    and dialects.  Hence there is more than one right answer to this
    question.}]

  For me (Francis), only (\ref{s:is-sg}) and (\ref{s:are-pl}) are ok;
  for many speakers (\ref{s:is-pl}) is also ok.

  
  
\item[B.] How would you elaborate or modify our analysis of the \lex{there}
  construction so as to capture the generalization you have discovered?
  Be as precise as you can.
\end{itemize}


\myslide{P2: Passing Up the Index}
\MyLogo{Based on \citet[Chapter 11, Problem 3]{sag:wasow:bender:2003}}

\begin{itemize}
\item[A.] Give the \ft{restr} value that our grammar should assign to
the sentence in~(i).  Be sure that the \ft{sit} value of the \prd{smoke}
predication
is identified with the \ft{annoyance} value of the \prd{annoy}
predication.

\begin{itemize}
\item[(i)] \eng{That Dana is smoking annoys Leslie.}
\end{itemize}

\noindent [{\sl Hint: This sentence involves two of the phenomena
  analyzed in this chapter: predicative complements of \lex{be} and CP subjects.}]

\item[B.] Draw a tree for (i).  Use abbreviations for node labels, but
  show the index on each node.

\item[C.] Explain how the \ft{sit} value of the \prd{smoke} 
%relation
predication
gets identified with the \ft{annoyance} value of the \prd{annoy} 
%relation.
predication.
Be sure to make reference to lexical entries, phrase structure rules,
and principles, as appropriate.

\end{itemize}


%\vbox{


\myslide{P3: An Annoying Problem}
\MyLogo{Based on \citet[Chapter 11, Problem 4]{sag:wasow:bender:2003}}

Assume that the lexical entry for the verb \lex{annoy} is the
following: 
%(Note that the first argument of this lexeme
%overrides the \index{defaults} default condition -- associated with the type {\it
%verb-lxm} -- requiring the first argument to be an NP.)

\begin{exe}
\ex\label{s:exi} {\avm{< \textnormal{annoy} ,\  [\type*{stv-lxm}\\
              arg-st & < [sem & [index & {\1}]] ,
                               NP$_i$ > \\
%			  [sem\ [index $i$]] > \\
              sem    & [index & $s$\\
                         restr  & < [reln & \prd{annoy}\\
                                       sit & \ \  $s$\\
                            	       annoyed & \ \ $i$\\
                            	       annoyance & \ \ {\1} ] > ]] >
}}
\end{exe}

\newpage
\begin{itemize}
% \item[A.] Show the lexeme that this lexical entry licenses, being sure
% to include all of the constraints it inherits from its supertypes.
% Show a lexical entry for {\it annoys} that results from
% applying the 3rd-Singular Verb Lexical Rule to this entry.
\item[A.] What constraints are imposed on the lexical sequences that
  result from applying the \txx{3rd-Singular Verb Lexical Rule} to
  this entry (including those that involve inheritance of constraints
  from the entry's supertypes)?
% \item[B.] Show the lexical entry that results from applying the
% \index{Extraposition Lexical Rule} Extraposition Lexical Rule to
% your answer to part (A).

\item[B.] What constraints are imposed on lexical sequences that
  result from applying the \txx{Extraposition Lexical Rule} to your
  answer to part (A)?

\item[C.] Draw a tree structure for the sentence in (\ref{s:exii}).
  You should show the value of all \ft{sem} features on all of the nodes,
  as well as the \ft{spr} and \ft{comps} features for \eng{annoys}.

\begin{exe}
  \ex\label{s:exii} \eng{It annoys Lee that Fido barks.}
\end{exe}
\newpage
\item[D.] The lexical entry for \lex{annoy} allows NP subjects as
  well, as in (\ref{s:exiii}).  Why doesn't the grammar then also
  license (\ref{s:exiv})?

\begin{exe}
\ex\label{s:exiii} \eng{Sandy annoys me.}

\ex\label{s:exiv} \bad{It annoys me Sandy.}
\end{exe}
\end{itemize}



\myslide{Acknowledgments and References}

\begin{itemize}
\item Course design and slides borrow heavily from Emily Bender's course:
\textit{Linguistics 566: Introduction to Syntax for Computational Linguistics}
\\ \url{http://courses.washington.edu/ling566}
\item This lecture complements \citet[Chapter 11]{sag:wasow:bender:2003}
\end{itemize}
\bibliographystyle{aclnat}
\bibliography{abb,mtg,nlp,ling}
\end{document}


%%% Local Variables: 
%%% coding: utf-8
%%% mode: latex
%%% TeX-PDF-mode: t
%%% TeX-engine: xetex
%%% End: 
